% GENERATED FILE. Edit canonical lessons, then run npm run generate:books.
% canonical-source-hash: fnv1a64:1d53290e3274f0e3
% canonical-lessons: TA-C498-ilainar, TA-C498-valipan, TA-C498-pirammaccari, TA-C498-varan, TA-C498-cammatam, TA-R498-review-492-498

\chapter{Young and Single}
\label{ch:ta-young-and-single}
\hlchaptermodality{498}

\begin{chapteropening}
\textbf{By the end of this chapter:} \emph{I can talk about young people, bachelors, finding a match and consent.}

Complete the last lesson of chapter 498: I can talk about young people, bachelors, finding a match and consent.
\end{chapteropening}

\section[\texorpdfstring{\ta{இளைஞர்} (iḷaiñar)}{இளைஞர் (iḷaiñar)}]{\ta{இளைஞர்} (iḷaiñar) — a young person}

\label{lesson:TA-C498-ilainar}



\begin{quote}
Before the new one: say the Tamil for to look after, then the Tamil for counsel.
\end{quote}

\subsection*{You'll want to know: \ta{இளைஞர்}}
\textbf{\ta{இளைஞர்}} — \emph{iḷaiñar} — \textquotedblleft{}a young person\textquotedblright{}.

\begin{quote}
\ta{விழாவில்} \ta{நிறைய} \ta{இளைஞர்கள்} \ta{கலந்துகொண்டார்கள்}.
\end{quote}

\begin{quote}
\emph{viḻāvil niṟaiya iḷaiñarkaḷ kalantukoṇṭārkaḷ.}
\end{quote}

\begin{quote}
Many young people took part in the festival.
\end{quote}

Plural: \ta{இளைஞர்கள்}, \emph{iḷaiñarkaḷ}.

Respectful -\ta{அர்} ending; plural \ta{இளைஞர்கள்}.

\begin{cousinweb}
\textbf{\ta{இளைஞர்}} (\emph{iḷaiñar}) is native Dravidian, built on \ta{இள} \emph{iḷa}, \textquotedblleft{}young, tender\textquotedblright{} (DEDR 513).
\end{cousinweb}

\begin{grammarlens}[title={What you've built}]
One more word for describing what you lived and felt.
\end{grammarlens}

\subsection*{Your turn}
\begin{itemize}
\raggedright
  \item \emph{Say it:} \emph{iḷaiñar}
  \item \emph{Say it:} \emph{iḷaiñar}, once more
  \item \emph{Say it:} \emph{viḻāvil niṟaiya iḷaiñarkaḷ kalantukoṇṭārkaḷ.} — Many young people took part in the festival.
  \item \emph{Recall:} say the Tamil for to look after, then the Tamil for counsel, then say \emph{iḷaiñar} again
\end{itemize}

\begin{tcolorbox}[breakable,colback=teal!4,colframe=teal!35!black,title={Before you move on}]
What does \ta{இளைஞர்} mean? (\textquotedblleft{}A young person\textquotedblright{}.) Say it once more.
\end{tcolorbox}
\section[\texorpdfstring{\ta{வாலிபன்} (vālipaṉ)}{வாலிபன் (vālipaṉ)}]{\ta{வாலிபன்} (vālipaṉ) — a young man}

\label{lesson:TA-C498-valipan}



\begin{quote}
Before the new one: say the Tamil for counsel, then the Tamil for a young person.
\end{quote}

\subsection*{You'll want to know: \ta{வாலிபன்}}
\textbf{\ta{வாலிபன்}} — \emph{vālipaṉ} — \textquotedblleft{}a young man\textquotedblright{}.

\begin{quote}
\ta{அந்த} \ta{வாலிபன்} \ta{தினமும்} \ta{காலையில்} \ta{ஓடுகிறான்}.
\end{quote}

\begin{quote}
\emph{anta vālipaṉ tiṉamum kālaiyil ōṭukiṟāṉ.}
\end{quote}

\begin{quote}
That young man runs every morning.
\end{quote}

Somewhat old-fashioned; \ta{இளைஞன்} is the neutral word.

\begin{cousinweb}
\textbf{\ta{வாலிபன்}} (\emph{vālipaṉ}) is usually traced to Sanskrit \emph{bāla}, \textquotedblleft{}young\textquotedblright{} (CDIAL 9216), with the Tamil -\ta{அன்} of a man.
\end{cousinweb}

\begin{grammarlens}[title={What you've built}]
One more word for describing what you lived and felt.
\end{grammarlens}

\subsection*{Your turn}
\begin{itemize}
\raggedright
  \item \emph{Say it:} \emph{vālipaṉ}
  \item \emph{Say it:} \emph{vālipaṉ}, once more
  \item \emph{Say it:} \emph{anta vālipaṉ tiṉamum kālaiyil ōṭukiṟāṉ.} — That young man runs every morning.
  \item \emph{Recall:} say the Tamil for counsel, then the Tamil for a young person, then say \emph{vālipaṉ} again
\end{itemize}

\begin{tcolorbox}[breakable,colback=teal!4,colframe=teal!35!black,title={Before you move on}]
What does \ta{வாலிபன்} mean? (\textquotedblleft{}A young man\textquotedblright{}.) Say it once more.
\end{tcolorbox}
\section[\texorpdfstring{\ta{பிரம்மச்சாரி} (pirammaccāri)}{பிரம்மச்சாரி (pirammaccāri)}]{\ta{பிரம்மச்சாரி} (pirammaccāri) — a bachelor}

\label{lesson:TA-C498-pirammaccari}



\begin{quote}
Before the new one: say the Tamil for a young person, then the Tamil for a young man.
\end{quote}

\subsection*{You'll want to know: \ta{பிரம்மச்சாரி}}
\textbf{\ta{பிரம்மச்சாரி}} — \emph{pirammaccāri} — \textquotedblleft{}a bachelor\textquotedblright{}.

\begin{quote}
\ta{என்} \ta{மாமா} \ta{இன்னும்} \ta{பிரம்மச்சாரியாக} \ta{இருக்கிறார்}.
\end{quote}

\begin{quote}
\emph{eṉ māmā iṉṉum pirammaccāriyāka irukkiṟār.}
\end{quote}

\begin{quote}
My uncle is still a bachelor.
\end{quote}

Also a celibate student in older usage.

\begin{cousinweb}
\textbf{\ta{பிரம்மச்சாரி}} (\emph{pirammaccāri}) is borrowed from Sanskrit \emph{brahmacārin}, \textquotedblleft{}an unmarried religious student\textquotedblright{} (Monier-Williams).
\end{cousinweb}

\begin{grammarlens}[title={What you've built}]
One more word for describing what you lived and felt.
\end{grammarlens}

\subsection*{Your turn}
\begin{itemize}
\raggedright
  \item \emph{Say it:} \emph{pirammaccāri}
  \item \emph{Say it:} \emph{pirammaccāri}, once more
  \item \emph{Say it:} \emph{eṉ māmā iṉṉum pirammaccāriyāka irukkiṟār.} — My uncle is still a bachelor.
  \item \emph{Recall:} say the Tamil for a young person, then the Tamil for a young man, then say \emph{pirammaccāri} again
\end{itemize}

\begin{tcolorbox}[breakable,colback=teal!4,colframe=teal!35!black,title={Before you move on}]
What does \ta{பிரம்மச்சாரி} mean? (\textquotedblleft{}A bachelor\textquotedblright{}.) Say it once more.
\end{tcolorbox}
\section[\texorpdfstring{\ta{வரன்} (varaṉ)}{வரன் (varaṉ)}]{\ta{வரன்} (varaṉ) — a prospective groom; a match}

\label{lesson:TA-C498-varan}



\begin{quote}
Before the new one: say the Tamil for a young man, then the Tamil for a bachelor.
\end{quote}

\subsection*{You'll want to know: \ta{வரன்}}
\textbf{\ta{வரன்}} — \emph{varaṉ} — \textquotedblleft{}a prospective groom; a match\textquotedblright{}.

\begin{quote}
\ta{மகளுக்கு} \ta{அப்பா} \ta{நல்ல} \ta{வரன்} \ta{தேடுகிறார்}.
\end{quote}

\begin{quote}
\emph{makaḷukku appā nalla varaṉ tēṭukiṟār.}
\end{quote}

\begin{quote}
Father is looking for a good match for his daughter.
\end{quote}

Said of an arranged marriage: \ta{வரன்} \ta{பார்}, to look for a match.

\begin{cousinweb}
\textbf{\ta{வரன்}} (\emph{varaṉ}) is borrowed from Sanskrit \emph{vara}, \textquotedblleft{}suitor, bridegroom\textquotedblright{} (CDIAL 11309).
\end{cousinweb}

\begin{grammarlens}[title={What you've built}]
One more word for describing what you lived and felt.
\end{grammarlens}

\subsection*{Your turn}
\begin{itemize}
\raggedright
  \item \emph{Say it:} \emph{varaṉ}
  \item \emph{Say it:} \emph{varaṉ}, once more
  \item \emph{Say it:} \emph{makaḷukku appā nalla varaṉ tēṭukiṟār.} — Father is looking for a good match for his daughter.
  \item \emph{Recall:} say the Tamil for a young man, then the Tamil for a bachelor, then say \emph{varaṉ} again
\end{itemize}

\begin{tcolorbox}[breakable,colback=teal!4,colframe=teal!35!black,title={Before you move on}]
What does \ta{வரன்} mean? (\textquotedblleft{}A prospective groom; a match\textquotedblright{}.) Say it once more.
\end{tcolorbox}
\section[\texorpdfstring{\ta{சம்மதம்} (cammatam)}{சம்மதம் (cammatam)}]{\ta{சம்மதம்} (cammatam) — consent}

\label{lesson:TA-C498-cammatam}



\begin{quote}
Before the new one: say the Tamil for a bachelor, then the Tamil for a prospective groom.
\end{quote}

\subsection*{You'll want to know: \ta{சம்மதம்}}
\textbf{\ta{சம்மதம்}} — \emph{cammatam} — \textquotedblleft{}consent\textquotedblright{}.

\begin{quote}
\ta{திருமணத்துக்குப்} \ta{பெண்ணின்} \ta{சம்மதம்} \ta{வேண்டும்}.
\end{quote}

\begin{quote}
\emph{tirumaṇattukkup peṇṇiṉ cammatam vēṇṭum.}
\end{quote}

\begin{quote}
A marriage needs the woman’s consent.
\end{quote}

\ta{சம்மதம்} \ta{சொல்}, to give consent.

\begin{cousinweb}
\textbf{\ta{சம்மதம்}} (\emph{cammatam}) is borrowed from Sanskrit \emph{sammata}, \textquotedblleft{}agreed, consented to\textquotedblright{} (Monier-Williams).
\end{cousinweb}

\begin{grammarlens}[title={What you've built}]
One more word for describing what you lived and felt.
\end{grammarlens}

\subsection*{Your turn}
\begin{itemize}
\raggedright
  \item \emph{Say it:} \emph{cammatam}
  \item \emph{Say it:} \emph{cammatam}, once more
  \item \emph{Say it:} \emph{tirumaṇattukkup peṇṇiṉ cammatam vēṇṭum.} — A marriage needs the woman’s consent.
  \item \emph{Recall:} say the Tamil for a bachelor, then the Tamil for a prospective groom, then say \emph{cammatam} again
\end{itemize}

\begin{tcolorbox}[breakable,colback=teal!4,colframe=teal!35!black,title={Before you move on}]
What does \ta{சம்மதம்} mean? (\textquotedblleft{}Consent\textquotedblright{}.) Say it once more.
\end{tcolorbox}
\section[(dialogue)]{Looking back: from Mother and Father to Young and Single}

\label{lesson:TA-R498-review-492-498}



\begin{quote}
Say the Tamil for one’s father and a mother.
\end{quote}

\subsection*{Your turn}
Say these aloud:

\begin{itemize}
\raggedright
  \item one’s father, a mother, a brother, a sister, the youngest child
  \item an elder brother’s wife, a husband’s younger brother, a husband’s sister, a co-parent-in-law, a mother’s younger sister
  \item a paternal kinsman, lineage, a generation, offspring, an heir
  \item an infant, a toddler, to crawl, to sing a lullaby, to feed
  \item pregnancy, childbirth, to be born, a son or daughter, to adopt
  \item to coo over, to pamper, to obey, to look after, counsel
  \item a young person, a young man, a bachelor, a prospective groom, consent
\end{itemize}

\begin{tcolorbox}[breakable,colback=teal!4,colframe=teal!35!black,title={Before you move on}]
One’s father? (\textbf{\ta{தந்தை}}, \emph{tantai}.) A mother? (\textbf{\ta{தாய்}}, \emph{tāy}.) A brother? (\textbf{\ta{சகோதரன்}}, \emph{cakōtaraṉ}.) A sister? (\textbf{\ta{சகோதரி}}, \emph{cakōtari}.) The youngest child? (\textbf{\ta{கடைக்குட்டி}}, \emph{kaṭaikkuṭṭi}.) An elder brother’s wife? (\textbf{\ta{அண்ணி}}, \emph{aṇṇi}.) A husband’s younger brother? (\textbf{\ta{கொழுந்தன்}}, \emph{koḻuntaṉ}.) A husband’s sister? (\textbf{\ta{நாத்தனார்}}, \emph{nāttaṉār}.) A co-parent-in-law? (\textbf{\ta{சம்பந்தி}}, \emph{campanti}.) A mother’s younger sister? (\textbf{\ta{சித்தி}}, \emph{citti}.) A paternal kinsman? (\textbf{\ta{பங்காளி}}, \emph{paṅkāḷi}.) Lineage? (\textbf{\ta{பரம்பரை}}, \emph{paramparai}.) A generation? (\textbf{\ta{தலைமுறை}}, \emph{talaimuṟai}.) Offspring? (\textbf{\ta{சந்ததி}}, \emph{cantati}.) An heir? (\textbf{\ta{வாரிசு}}, \emph{vāricu}.) An infant? (\textbf{\ta{சிசு}}, \emph{cicu}.) A toddler? (\textbf{\ta{மழலை}}, \emph{maḻalai}.) To crawl? (\textbf{\ta{தவழ்}}, \emph{tavaḻ}.) To sing a lullaby? (\textbf{\ta{தாலாட்டு}}, \emph{tālāṭṭu}.) To feed? (\textbf{\ta{ஊட்டு}}, \emph{ūṭṭu}.) Pregnancy? (\textbf{\ta{கர்ப்பம்}}, \emph{karppam}.) Childbirth? (\textbf{\ta{பிரசவம்}}, \emph{piracavam}.) To be born? (\textbf{\ta{பிற}}, \emph{piṟa}.) A son or daughter? (\textbf{\ta{பிள்ளை}}, \emph{piḷḷai}.) To adopt? (\textbf{\ta{தத்தெடு}}, \emph{tatteṭu}.) To coo over? (\textbf{\ta{கொஞ்சு}}, \emph{koñcu}.) To pamper? (\textbf{\ta{சீராட்டு}}, \emph{cīrāṭṭu}.) To obey? (\textbf{\ta{கீழ்ப்படி}}, \emph{kīḻppaṭi}.) To look after? (\textbf{\ta{பராமரி}}, \emph{parāmari}.) Counsel? (\textbf{\ta{ஆலோசனை}}, \emph{ālōcaṉai}.) A young person? (\textbf{\ta{இளைஞர்}}, \emph{iḷaiñar}.) A young man? (\textbf{\ta{வாலிபன்}}, \emph{vālipaṉ}.) A bachelor? (\textbf{\ta{பிரம்மச்சாரி}}, \emph{pirammaccāri}.) A prospective groom? (\textbf{\ta{வரன்}}, \emph{varaṉ}.) Consent? (\textbf{\ta{சம்மதம்}}, \emph{cammatam}.)
\end{tcolorbox}
